\documentclass{article}
\usepackage[utf8]{inputenc}
\usepackage{amsmath}

\begin{document}

In this study, the transient gas flow equation for the microporous region is solved using **parameter explicit lagging** combined with **semi-implicit temporal discretization**. To ensure convergence and accuracy, a **Picard iteration** scheme is employed for updating variables when necessary. The governing equation is expressed as:

\begin{equation}
    \frac{\partial}{\partial t}\left[\left(\varphi-\frac{n_{ad}}{\rho^{ad}}\right) \rho^{g}+n_{ad}\right]+\nabla \cdot\left(\rho^{g} \mathbf{q}^{g}\right)+\nabla \cdot\left(\rho^{ad} \mathbf{v}^{ad}\right)=0
\end{equation}

To solve this, we treat the storage and transport terms separately. First, regarding the storage term, let:

\begin{equation}
    S(p)=\left(\varphi -\frac{n_{ad}}{\rho^{ad}}\right) \rho^{g}+n_{ad}
\end{equation}

Applying the chain rule, the temporal derivative of the storage term is given by:

\begin{equation}
    \frac{\partial S(p)}{\partial t}=\frac{d S(p)}{d p} \cdot \frac{\partial p}{\partial t}
\end{equation}

The derivative of $S(p)$ with respect to pressure $p$ is expanded as:

\begin{equation}
    \frac{d S(p)}{d p}=\left(\varphi-\frac{n_{ad}}{\rho^{ad}}\right) \frac{d \rho^{g}}{d p}+\rho^{g}\left(-\frac{1}{\rho^{ad}} \frac{d n_{ad}}{d p}\right)+\frac{d n_{ad}}{d p}
\end{equation}

Rearranging Equation 4 yields:

\begin{equation}
    \frac{d S(p)}{d p}=\left(\phi-\frac{n_{ad}}{\rho^{ad}}\right) \frac{d \rho^{g}}{d p}+\left(1-\frac{\rho^{g}}{\rho^{ad}}\right) \frac{d n_{ad}}{d p}
\end{equation}

By introducing the isothermal compressibility coefficient for real gases, $c_g$:

\begin{equation}
    c_{g}=\frac{1}{\rho^{g}} \frac{d \rho^{g}}{d p}=\frac{1}{p}-\frac{1}{Z} \frac{d Z}{d p}
\end{equation}

Since the pressure change is minimal, the compressibility factor can be assumed to be approximately constant, i.e., $\frac{dZ}{dp} = 0$. Consequently, Equation 6 simplifies to:

\begin{equation}
    c_{g}=\frac{1}{p}
\end{equation}

The derivatives of gas density and adsorption capacity with respect to pressure are obtained as follows:

\begin{equation}
    \frac{d \rho^{g}}{d p}=\rho^{g} c_{g}=\frac{M p}{Z R T} c_{g}=\frac{M}{Z R T}
\end{equation}

\begin{equation}
    \frac{d n_{a d}}{d p}=\frac{n_{a d}^{\max } K}{(1+K p)^{2}}
\end{equation}

We can then define a **total storage capacity** ($C$), where $C(p) = \frac{dS(p)}{dp}$. Substituting the previous expressions, we obtain:

\begin{equation}
    C_{m i c r o}(p)=\left(\varphi-\frac{n_{a d}^{\max } K p}{\rho^{a d}(1+K p)}\right) \frac{M}{Z R T}+\left(1-\frac{M p}{Z R T \rho^{a d}}\right) \frac{n_{a d}^{\max } K}{(1+K p)^{2}}
\end{equation}

In the transport terms, nonlinearities are introduced by the compressibility of the free gas and the surface diffusion of the adsorbed gas:

\begin{equation}
    \mathbf{q}^{g}=-\text{fun}(\text{Kn}) \frac{k_{m}}{\mu} \nabla p
\end{equation}

\begin{equation}
    \mathbf{v}^{a d}=-\frac{1}{\rho^{a d}} D_{s} \nabla n_{a d}
\end{equation}

\begin{equation}
    \nabla \cdot\left(\rho^{g} \mathbf{q}^{g}\right)+\nabla \cdot\left(\rho^{a d} \mathbf{v}^{a d}\right)=\nabla \cdot\left(\rho^{g} \text{fun}(\text{Kn}) \frac{k_{m}}{\mu}+\frac{K n_{a d}^{\max } D_{s}}{(1+K p)^{2}}\right) \nabla p
\end{equation}

By defining the **total transmissibility coefficient** as:

\begin{equation}
    T_{m i c r o}(p)=\left(\rho^{g} \text{fun}(\text{Kn}) \frac{k_{m}}{\mu}+\frac{K n_{a d}^{\max } D_{s}}{(1+K p)^{2}}\right)
\end{equation}

By substituting the total storage capacity and the total transmissibility coefficient into the original governing equation, the updated governing equation for gas flow in the micropores is obtained:

\begin{equation}
    C_{m i c r o}(p) \cdot \frac{\partial p}{\partial t}=\nabla \cdot\left(T_{m i c r o}(p) \nabla p\right)
\end{equation}

Although the equation has been simplified in form, the coefficient $C(p)$ still contains highly nonlinear terms related to pressure $p$. To achieve a linearized solution, these coefficients must be "frozen" as constants at a specific reference pressure $p^{b}$. In this study, we take the pressure from background pressure, $p^b$, as the reference pressure. Consequently, the final governing equation for gas flow becomes a linear equation with respect to pressure:

\begin{equation}
    C_{m i c r o}(p^{b}) \cdot \frac{\partial p}{\partial t}=\nabla \cdot\left(T_{m i c r o}(p^{b}) \nabla p\right)
\end{equation}

For gas flow in the macropores, only the free gas phase is considered:

\begin{equation}
    V_{i} \frac{\partial \rho_{i}^{g}}{\partial t}+\sum_{j=1}^{N_{i}} T_{i j}\left(p_{i}-p_{j}\right)=0
\end{equation}

Similarly, we define the storage capacity for the macropores as $C_{\text{macro}}(p)$:

\begin{equation}
    C_{\text{macro}}(p)=\frac{M}{Z(p) R T}
\end{equation}

By defining the reference pressure $p^{b}$ as the pressure from the background pressure, Equation 17 can be rewritten as:

\begin{equation}
    V_{i} C_{i, \text{macro}}(p^{b}) \frac{\partial p_{i}}{\partial t}+\sum_{j=1}^{N_{i}} T_{i j}\left(p_{i}-p_{j}\right)=0
\end{equation}

Equation 16 can then be discretized as follows:

\begin{equation}
    V_{i} C_{i, \text{micro}}\left(p^{b}\right) \frac{p_{i}^{t+\Delta t}-p_i^{t}}{\Delta t}+\sum_{j=1}^{N_{i}} T_{i j, \text{micro}}\left(p^{b}\right)\left(p_{i}^{t+\Delta t}-p_{j}^{t+\Delta t}\right)=0
\end{equation}

The transmissibility coefficient between micropores, $T_{i j,micro}\left(p^{b}\right)$, is given by:

\begin{equation}
    T_{i j, \text{micro}}\left(p^{b}\right)=\frac{a_{i} a_{j}}{a_{i}+a_{j}}, \text{ with } a_{i}=\frac{A_{i} T_{i, \text{micro}}\left(p^{b}\right)}{d_{i}} \mathbf{n}_{i} \cdot \mathbf{f}_{i} \text{ and } a_{j}=\frac{A_{j} T_{j, \text{micro}}\left(p^{b}\right)}{d_{j}} \mathbf{n}_{j} \cdot \mathbf{f}_{j}
\end{equation}

Equation 19 is discretized as:

\begin{equation}
    V_{i} C_{i, \text{macro}}\left(p^{b}\right) \frac{p_{i}^{t+\Delta t}-p_i^{t}}{\Delta t}+\sum_{j=1}^{N_{i}} T_{i j, \text{macro}}\left(p^{b}\right)\left(p_{i}^{t+\Delta t}-p_{j}^{t+\Delta t}\right)=0
\end{equation}

The gas transmissibility coefficient  between macropores $T_{i j, \text{macro}}\left(p^{b}\right)$ is defined as:

\begin{equation}
    T_{i j, \text{macro}}\left(p^{b}\right)=\frac{\operatorname{fun}_{i}(\mathrm{Kn})+\operatorname{fun}_{j}(\mathrm{Kn})}{2} \frac{\pi R_{i j}^{4}}{8 \mu l_{i j}} * \frac{\rho_{i}^{t}+\rho_{j}^{t}}{2}
\end{equation}

Finally, the gas transmissibility coefficient between the macropores and micropores is defined as:

\begin{equation}
    T_{i j}\left(p^{b}\right)=\frac{a_{i} a_{j}}{a_{i}+a_{j}}
\end{equation}

where the interface terms $a_i$ and $a_j$ are given by:

\begin{equation}
    a_{i}=\rho_{i}^{t} \operatorname{fun}_{i}(\mathrm{Kn}) \frac{\pi R_{i}^{4}}{8 \mu l_{i}} \quad \text{and} \quad a_{j}=\frac{A_{j} T_{j, micro}\left(p^{b}\right)}{d_{j}} \mathbf{n}_{j} \cdot \mathbf{f}_{j}
\end{equation}

\end{document}